% $Header: /Users/joseph/Documents/LaTeX/beamer/solutions/conference-talks/conference-ornate-20min.en.tex,v 90e850259b8b 2007/01/28 20:48:30 tantau $

\documentclass[14pt,t]{beamer}
\usepackage[fontset=ubuntu]{ctex}
\usepackage[english]{babel}
\geometry{papersize={160mm,90mm}}
%\usepackage{times}
%\usepackage[T1]{fontenc}


\mode<presentation>
{%
  \usetheme{Warsaw}%
  \setbeamercovered{transparent}%
}

\title[树莓派体验] % (optional, use only with long paper titles)
{树莓派体验}

\subtitle{}

% (optional, use only with lots of authors)
\author[Fang Yuan (yfang@nju.edu.cn)]
{方元 (yfang@nju.edu.cn)}

\institute[南京大学电子学院] % (optional, but mostly needed)
{南京大学电子学院}

\date[expcourse] % (optional, should be abbreviation of conference name)
{\today}

\subject{电子计算机研究与应用}
% This is only inserted into the PDF information catalog. Can be left
% out. 

% If you have a file called "university-logo-filename.xxx", where xxx
% is a graphic format that can be processed by latex or pdflatex,
% resp., then you can add a logo as follows:

\pgfdeclareimage[height=7mm]{university-logo}{njulogo.pdf}
\logo{\pgfuseimage{university-logo}}

% Delete this, if you do not want the table of contents to pop up at
% the beginning of each subsection:
%\AtBeginSubsection[]%
%{%
%  \begin{frame}<beamer>{Agenda}
%    \tableofcontents[currentsection,currentsubsection]
%  \end{frame}
%}

% If you wish to uncover everything in a step-wise fashion, uncomment
% the following command: 

%\beamerdefaultoverlayspecification{<+->}

\begin{document}

\begin{frame}
  \titlepage
\end{frame}

\begin{frame}{Agenda}
    \small
  \tableofcontents
  % You might wish to add the option [pausesections]
\end{frame}


% Structuring a talk is a difficult task and the following structure
% may not be suitable. Here are some rules that apply for this
% solution: 

% - Exactly two or three sections (other than the summary).
% - At *most* three subsections per section.
% - Talk about 30s to 2min per frame. So there should be between about
%   15 and 30 frames, all told.

% - A conference audience is likely to know very little of what you
%   are going to talk about. So *simplify*!
% - In a 20min talk, getting the main ideas across is hard
%   enough. Leave out details, even if it means being less precise than
%   you think necessary.
% - If you omit details that are vital to the proof/implementation,
%   just say so once. Everybody will be happy with that.

\section{计算机操作系统}

\subsection{什么是计算机}

\begin{frame}{什么是计算机}

\begin{figure}
\centering
\includegraphics[width=.5\textwidth]{computers.png}
\end{figure}
\end{frame}

\begin{frame}{什么是计算机}

\begin{center}
\includegraphics[width=.2\textwidth]{handphone.jpg}
    \hspace{2mm}
\includegraphics[width=.2\textwidth]{wearables.jpg}
    \hspace{2mm}
\includegraphics[width=.3\textwidth]{router.jpg}
\end{center}

\end{frame}

\subsection{操作系统}

\begin{frame}{操作系统}
    操作系统是计算机的核心软件, 它控制和管理整个计算机系统硬件和软件资源，
    以提供给用户和其它软件方便的接口和环境。

\begin{center}
\includegraphics[width=.2\textwidth]{windowslogo.jpg}
\end{center}

\end{frame}

\begin{frame}{桌面操作系统}

    \ \\
\begin{center}
\includegraphics[width=.15\textwidth]{android.png}
    \hspace{3mm}
\includegraphics[width=.15\textwidth]{macos.pdf}
    \hspace{3mm}
\includegraphics[width=.15\textwidth]{chrome.pdf}
    \hspace{3mm}
\includegraphics[width=.15\textwidth]{apple.jpg}
\end{center}

\end{frame}

\subsection{Linux}

\begin{frame}{Linus 和 Linux}

\begin{center}
\includegraphics[width=.2\textwidth]{Linus.jpg}
    \hspace{2cm}
\includegraphics[width=.2\textwidth]{tux.pdf}
\end{center}
    1991 年，芬兰大学生Linus(李纳斯)在学习计算机课程期间，
    出于兴趣开发了 Linux 操作系统。
\end{frame}

\begin{frame}{Linix 发行版}
\begin{center}
\includegraphics[width=.5\textwidth]{linux_distros.pdf}
\end{center}
\end{frame}

\begin{frame}{Linux 用在哪里}
\begin{itemize}
    \item 服务器、超级计算机 (Top 500 全部采用Linux)
    \item 手机、平板 (Linux 内核约占 70\% 的市场份额)
    \item 机顶盒、网络路由器、可穿戴设备...
    \item 个人电脑
\end{itemize}
\end{frame}

\begin{frame}{Linux 的特点}
\begin{itemize}
    \item 自由软件 (free software)，遵循开源版权协议
    \item 标准化
    \item 大量的软件支持
    \item 安全、可靠
    \item 网络性能优越
\end{itemize}
\end{frame}

\section{Python}

\subsection{编程语言}
\begin{frame}{编程语言}
    编程语言是一种计算机和人都能识别的符号体系。
    一种计算机语言让程序员能够准确地定义计算机所需要使用的数据，
    并精确地定义在不同情况下应当产生的动作。
\end{frame}

\begin{frame}{编程语言}
    计算机语言从低级 (接近机器)到高级 (接近人类语言)可分为
\begin{itemize}
    \item 机器语言 (0、1编码)
    \item 汇编语言 (简单的符号)
    \item 高级语言 (完善的符号)
    \item 自然语言 (发展中)
\end{itemize}
所有编程语言都需要转换到机器语言才能让计算机执行。
编译型/解释型语言通过不同的方式转换。
\end{frame}

\subsection{Python}
\begin{frame}{Python}
\begin{itemize}
    \item 高级、通用、解释型编程语言
    \item 语法简洁, 用更少的代码表达设计思想，强调可读性
    \item 大量的第三方库支持
\end{itemize}
    Python 语言有大量的支持库，应用于数值计算、图形处理、
    数据库、互联网、人工智能等几乎所有的计算领域。
\end{frame}

\begin{frame}{Python之父}
\begin{center}
\includegraphics[width=.16\textwidth]{Guido.jpg}
    \hspace{2cm}
\includegraphics[width=.2\textwidth]{python.pdf}
\end{center}

    Guido van Rossum (罗素姆) 荷兰程序员, 1989年圣诞节期间开发了
    Python 程序设计语言。
\end{frame}

\section{树莓派}
\subsection{树莓派}

\begin{frame}{树莓派的特点}
    树莓派是由英国的一个非盈利组织开发的用于中小学计算机教育的
    单板计算机。

    \begin{center}
        \includegraphics[width=.3\textwidth]{Pi1_A.jpg}
    \end{center}
\end{frame}

\begin{frame}{树莓派的特点}
\begin{itemize}
    \item 性价比
    \item 丰富的接口，可用于二次开发
    \item 支持高级操作系统
\end{itemize}
大多数树莓派应用基于Linux操作系统，Linux操作系统支持 Python 编程语言，
这使得基于树莓派的应用开发相当简单。
\end{frame}

\begin{frame}{在树莓派上编程}
    程序的层次结构
\begin{itemize}
    \item 常量、变量
    \item 表达式
    \item 语句
    \item 函数
    \item 库 (模块)
    \item 程序 (工程)
\end{itemize}
\end{frame}
\end{document}
